\section{Discussion}

\subsection{Why Grounding Did Not Move Factuality}

We set out expecting the knowledge graph to reduce unsupported claims and it
does not, at crop level, by any margin we can measure. The reason is not that
the graph fails but that the question was posed at the wrong unit. A crop-level
ladder gives every condition the same inventory for the crop under description
and then scores only claims about that crop; the retrieved history is about
other crops, so it has almost nothing to act on. Reported as an ablation it
would be noise presented as a result. At neighbourhood level, where the
conditions become genuinely different representations of the same observations,
the graph earns a claim --- and it is a counting claim, not a factuality one.
We state it that way rather than the way we intended to.

A second null belongs here for the same reason. On our data a deterministic
template is indistinguishable from a language model on CFS, and at
neighbourhood level it leads on CFS-F1. The case for generation rests on
readability and on producing a count a reader can act on; we measure only the
second.

\subsection{Limitations}

\textbf{We do not approach the supervised ceiling.} End-to-end models trained
on these benchmarks' splits reach 86--91 F1 on LEVIR-CD; we reach
\EOnePixelFOne. AnyChange is included as zero-shot context, not as a
supervision-matched baseline, because our small pairing head uses LEVIR-CD
training masks.

\textbf{The LEVIR-CD result is not zero-shot.} The detector never saw the
dataset, but the head was trained on its training split from mask-derived
labels. The WHU-CD transfer carries the zero-shot claim.

\textbf{The temporal half of the graph is untested.} The graph is designed to
accumulate entity histories across many revisits, but both benchmarks carry
exactly two timepoints, so what we exercise is spatial aggregation over the
crops of one tile, not history over time. The design's main hypothesis is
therefore not evaluated here.

\textbf{Claim extraction is rule-based and not independently validated.} This
is what keeps a language model out of the metric, but the extractor sees only
the direction cues, object classes and quantifiers it encodes, and is tuned to
this domain's vocabulary. Agreement with independently annotated claims should
be established before treating CFS as a general report-factuality metric.

\textbf{CFS lacks an independently generated specialist baseline.} Published
change-captioning outputs were unavailable in a form that we could rescore, so
all CFS rows use outputs generated in our evaluation harness. Conventional
caption metrics provide published context, but do not validate the new metric.

\textbf{Generation uncertainty is not quantified.} Each language-model
condition is represented by one completed generation pass. The small
crop-level differences among structured, flat-retrieval and graph conditions
are therefore descriptive rather than statistically significant effects.

\textbf{Both benchmarks are building-centric and aerial rather than
satellite}, so generalisation to other change types and lower ground sampling
distances is untested.

\textbf{The report model is licensed for research use only.} \LLMName{} is
released under a non-commercial licence, so the service this paper is written
around would swap in a commercially licensed model. Nothing in the pipeline
depends on which model writes the report, but nothing here shows the factuality
numbers survive that swap either.

\subsection{Future Work}

The experiment this work now needs is a dataset with a real revisit history:
more than two acquisitions of the same area, which is what would test the
temporal aggregation the graph was built for and which neither public benchmark
provides. Beyond that, change types other than buildings, true satellite ground
sampling distances, and a commercially licensed report model are the three
steps between this evaluation and a service that could ship. Independent manual
validation of the claim extractor and paired multi-seed generation tests are
the immediate evaluation priorities.
